\section{Булевы функции}
\label{sec:boolean-functions}

\subsection{Основные определения и число булевых функций}

Булевой функцией от $n$ переменных называется отображение
\[
    f\colon \{0,1\}^{n}\to\{0,1\}.
\]
Аргумент функции --- булев вектор $\alpha=(\alpha_1,\ldots,\alpha_n)$.
Так как имеется $2^n$ различных наборов аргументов и на каждом из них значение
функции выбирается независимо из двух значений, число всех $n$-местных булевых
функций равно
\[
    2^{2^n}.
\]

Переменная $x_i$ называется \emph{существенной} для $f$, если существуют два
набора, отличающиеся только значением $x_i$, на которых функция принимает разные
значения. В противном случае $x_i$ называется фиктивной.

К базовым операциям относятся
\[
 \neg x,\qquad x\land y,\qquad x\lor y,\qquad x\oplus y,
 \qquad x\to y=\neg x\lor y.
\]
Кроме того, важны штрих Шеффера и стрелка Пирса:
\[
 x\mid y=\neg(x\land y),\qquad
 x\downarrow y=\neg(x\lor y).
\]
Обе последние операции по отдельности образуют функционально полный базис.

Основные тождества булевой алгебры:
\[
 x\land x=x,\quad x\lor x=x,\quad
 x\land 1=x,\quad x\lor0=x,
\]
\[
 x\land0=0,\quad x\lor1=1,\quad
 x\land\neg x=0,\quad x\lor\neg x=1,
\]
а также коммутативность, ассоциативность, дистрибутивность и законы де Моргана
\[
 \neg(x\land y)=\neg x\lor\neg y,
 \qquad
 \neg(x\lor y)=\neg x\land\neg y.
\]

\subsection{Суперпозиция и разложение Шеннона}

Если $g(y_1,\ldots,y_m)$ и $f_1,\ldots,f_m$ --- булевы функции, то функция
\[
 h(x_1,\ldots,x_n)=g\bigl(f_1(x),\ldots,f_m(x)\bigr)
\]
называется их суперпозицией. Именно замыкание относительно суперпозиции лежит в
основе понятия функциональной полноты.

Для любой переменной $x_i$ справедливо \textbf{разложение Шеннона}
\[
 f(x_1,\ldots,x_n)=
 \bigl(\neg x_i\land f|_{x_i=0}\bigr)
 \lor
 \bigl(x_i\land f|_{x_i=1}\bigr).
\]
Проверка элементарна: при $x_i=0$ остаётся только первое слагаемое, при $x_i=1$
--- только второе. Последовательно раскладывая функцию по всем переменным,
получают канонические нормальные формы.

Разложение Шеннона полезно не только теоретически. Оно лежит в основе
декомпозиции логических схем и двоичных диаграмм решений (BDD): вершина диаграммы
соответствует выбору значения некоторой переменной, а две исходящие ветви ---
подстановкам $x_i=0$ и $x_i=1$.

\subsection{ДНФ, КНФ и совершенные нормальные формы}

\emph{Литералом} называется переменная или её отрицание. Конъюнкция литералов
называется элементарной конъюнкцией, дизъюнкция литералов --- элементарной
дизъюнкцией.

Дизъюнктивная нормальная форма (ДНФ) --- дизъюнкция элементарных конъюнкций,
а конъюнктивная нормальная форма (КНФ) --- конъюнкция элементарных дизъюнкций.

Минтерм для набора $\alpha\in\{0,1\}^n$ имеет вид
\[
 K_\alpha=x_1^{\alpha_1}\land\cdots\land x_n^{\alpha_n},
 \qquad
 x_i^1=x_i,\quad x_i^0=\neg x_i,
\]
и равен единице ровно на наборе $\alpha$. Поэтому для нетождественно нулевой
функции
\[
 f=\bigvee_{\alpha:\,f(\alpha)=1}K_\alpha
\]
--- её совершенная ДНФ (СДНФ).

Аналогично, для каждого нулевого набора строится макстерм, равный нулю только на
этом наборе; конъюнкция таких макстермов образует СКНФ. Для нетривиальной функции
СДНФ и СКНФ единственны с точностью до перестановки членов. Для константных
функций удобно считать пустую дизъюнкцию равной $0$, а пустую конъюнкцию --- $1$.

\subsection{Полином Жегалкина}

Над полем $\mathbb F_2=\{0,1\}$ сложению соответствует операция XOR, а
умножению --- конъюнкция. Любая булева функция единственным образом имеет вид
\[
 f(x_1,\ldots,x_n)=
 \bigoplus_{S\subseteq\{1,\ldots,n\}}
 a_S\prod_{i\in S}x_i,
 \qquad a_S\in\mathbb F_2.
\]
Это \emph{полином Жегалкина}, или алгебраическая нормальная форма.

Единственность можно объяснить двумя способами.
Во-первых, имеется ровно $2^n$ мономов
\[
 1,\ x_i,\ x_ix_j,\ldots,x_1\cdots x_n,
\]
и коэффициенты $a_S$ образуют вектор из $2^n$ бит; следовательно, таких полиномов
ровно $2^{2^n}$, то есть столько же, сколько булевых функций. Остаётся заметить,
что два различных приведённых полинома не задают одну и ту же функцию.
Во-вторых, коэффициенты можно последовательно восстановить по значениям функции:
$a_\varnothing=f(0,\ldots,0)$, затем коэффициенты мономов первой степени, второй
и так далее.

Функция называется линейной (точнее, аффинной над $\mathbb F_2$), если степень
полинома Жегалкина не превосходит единицы:
\[
 f=a_0\oplus a_1x_1\oplus\cdots\oplus a_nx_n.
\]

\subsection{Двойственность, монотонность и самодвойственность}

Двойственная функция определяется формулой
\[
 f^*(x_1,\ldots,x_n)=\neg f(\neg x_1,\ldots,\neg x_n).
\]
Если некоторое тождество булевой алгебры верно, то после взаимной замены
$\land\leftrightarrow\lor$ и $0\leftrightarrow1$ получается также верное
двойственное тождество --- принцип двойственности.

Функция самодвойственна, если $f=f^*$. В частности, самодвойственная функция не
может одновременно сохранять и $0$, и $1$ произвольным образом: её значения на
противоположных наборах $\alpha$ и $\neg\alpha$ связаны равенством
\[
 f(\neg\alpha)=\neg f(\alpha).
\]

На $\{0,1\}^n$ вводится покоординатный порядок:
$\alpha\le\beta$, если $\alpha_i\le\beta_i$ для всех $i$. Функция монотонна,
если
\[
 \alpha\le\beta\quad\Longrightarrow\quad f(\alpha)\le f(\beta).
\]

\subsection{Функциональная полнота и теорема Поста}

Для системы функций $F$ обозначим через $[F]$ множество всех функций, которые
можно получить из $F$ суперпозициями и подстановками переменных. Система
называется \emph{функционально полной}, если $[F]$ содержит все булевы функции.

В критерии Поста возникают пять важнейших замкнутых классов:
\begin{enumerate}
 \item $T_0$: $f(0,\ldots,0)=0$;
 \item $T_1$: $f(1,\ldots,1)=1$;
 \item $S$: самодвойственные функции;
 \item $M$: монотонные функции;
 \item $L$: линейные (аффинные над $\mathbb F_2$) функции.
\end{enumerate}
Каждый из этих классов замкнут относительно суперпозиции. Например, суперпозиция
функций, сохраняющих ноль, снова сохраняет ноль; суперпозиция монотонных функций
монотонна и т.~д. Поэтому если вся система $F$ лежит хотя бы в одном таком
классе, полной она быть не может.

\textbf{Теорема Поста.}
Система булевых функций $F$ функционально полна тогда и только тогда, когда она
не содержится целиком ни в одном из классов
\[
 T_0,\quad T_1,\quad S,\quad M,\quad L.
\]
Иными словами, в $F$ должны существовать (не обязательно одна и та же) функции,
нарушающие каждое из пяти свойств.

Практическая проверка полноты сводится к таблице: для каждой функции системы
проверяют принадлежность пяти классам. Если в каждом столбце найдётся хотя бы
одна функция, не принадлежащая соответствующему классу, система полна.

Например, один штрих Шеффера образует полный базис, поскольку
\[
 \neg x=x\mid x,
 \qquad
 x\land y=(x\mid y)\mid(x\mid y),
\]
после чего через $\neg$ и $\land$ можно построить любую СДНФ.

\subsection{Импликанты и минимизация ДНФ}

Элементарная конъюнкция $K$ называется \emph{импликантой} функции $f$, если
$K=1$ влечёт $f=1$. Импликанта называется простой, если удаление любого литерала
из неё разрушает это свойство. Сокращённая ДНФ строится из простых импликант.

Различают несколько критериев минимальности: минимальное число конъюнкций,
минимальное число литералов, минимальную аппаратную стоимость. Поэтому
``минимальная ДНФ'' всегда понимается относительно выбранного критерия.

Основные практические методы:
\begin{itemize}
 \item карты Карно --- геометрическое объединение соседних единичных клеток в
       группы размеров $1,2,4,\ldots$;
 \item метод Квайна--Мак-Класки --- систематическое объединение минтермов,
       построение простых импликант и таблицы покрытий;
 \item алгебраические преобразования с законами поглощения и склеивания.
\end{itemize}

\subsection{Что важно сказать на экзамене}

Нужно дать определение булевой функции, показать канонические представления
(СДНФ/СКНФ и полином Жегалкина) и объяснить суперпозицию. Центральный теоретический
результат --- критерий Поста: система полна тогда и только тогда, когда она не
лежит целиком ни в одном из пяти замкнутых классов $T_0,T_1,S,M,L$. Полезно
уметь привести полный базис, например штрих Шеффера, и коротко объяснить
минимизацию ДНФ через простые импликанты.

\newpage
